\chapter{Testing, Reproducibility and Performance}\label{ch:testing}

\section{Testing strategy}
The backend has about 115 files in \code{backend/\allowbreak tests} (881 test items collected by \code{pytest}, 823 \code{def test\_} functions by a plain text count; parametrised tests account for the difference). The worker package has nine tests and the frontend has 28 test files. Continuous integration (\code{.github/\allowbreak workflows/\allowbreak ci.yml}) runs on pushes to \code{Main} and on pull requests: backend \code{ruff} and \code{pytest} against a Postgres service container, worker \code{ruff} and \code{pytest} with a slim install, and the frontend tests, type-check and build. We did not generate a coverage report for this thesis, so we give no coverage percentage.

\paragraph{Latest full run on the reference machine.}
\begin{table}[H]
\centering\small
\caption{Latest test runs (2026-10-07, working tree on branch \code{refactor/\allowbreak remove-dead-code}).}
\begin{tabularx}{\textwidth}{@{}p{3.0cm}L@{}}
\toprule
\textbf{Suite} & \textbf{Result}\\
\midrule
Backend \code{pytest} & 661 passed, 220 skipped, 2 warnings (about 32--94\,s). Of the 220 skips, 216 are database-backed tests that skip here because no Postgres server was running on the reference machine (``connection refused''), one needs WeasyPrint system libraries, one is a live Hub import that needs \code{TRUSTLENS\_LIVE\_TESTS=1}. These DB tests run in CI, which provides Postgres; we did not run them in this session, so their status is not part of this table.\\
Worker \code{pytest} & 9 passed.\\
Frontend \code{vitest} & 28 files, 273 tests passed.\\
\code{ruff} (backend) & 17 findings. The baseline at the start of round~3 was 31.\\
\bottomrule
\end{tabularx}
\end{table}

\begin{table}[H]
\centering\small
\caption{Groups of automated tests (file names from \code{backend/\allowbreak tests}).}
\begin{tabularx}{\textwidth}{@{}p{3.4cm}L@{}}
\toprule
\textbf{Area} & \textbf{Examples and what they check}\\
\midrule
Scoring & \code{test\_fries\_scorer}: cube-root, veto, aspect mean, weights; compares against frozen vectors in \code{shared/\allowbreak scoring/\allowbreak fixtures/\allowbreak fries\_test\_vectors.json}.\\
Fairness & \code{test\_fairness\_stats}, \code{test\_fairness\_binary\_risk} (decision order of Table~\ref{tab:fairrule}), \code{test\_fairness\_base\_rate}, \code{test\_fairness\_confirm\_setup}, \code{test\_run\_fairness\_confirm}.\\
Robustness and safety & \code{test\_robustness\_eval}, \code{test\_robustness\_nlp}, \code{test\_prediction\_gate}, \code{test\_safety\_behavior}, \code{test\_safety\_probe\_behavior}, \code{test\_defect\_severity}, \code{test\_severity\_sweep}.\\
Integrity and cards & \code{test\_artifact\_integrity} (sharded sets, manifests, statuses), \code{test\_tamper\_demo}, \code{test\_integrity\_probe}, \code{test\_explainability\_card}, \code{test\_disclosure\_gaps\_vs\_risks} (L3 split).\\
O/S/D engines & \code{test\_deterministic\_osd}, \code{test\_osd\_agent}, \code{test\_osd\_calibration} (frozen anchors, \code{check\_frozen}), \code{test\_hybrid\_osd\_agent} (fallback), \code{test\_human\_review}.\\
Validation harnesses & \code{test\_l2\_interpretation}, \code{test\_l10\_validation}, \code{test\_claude\_eval}: packet building, leakage checks (keys and strings), hash matching, parsing, agreement categories.\\
Pipeline and API & \code{test\_evaluation\_enqueue}, \code{test\_pipeline\_missing\_evaluation}, \code{test\_evaluation\_lifecycle}, \code{test\_api\_reports}, \code{test\_mode\_disclosure\_scoring\_state}.\\
\bottomrule
\end{tabularx}
\end{table}

\paragraph{Test-driven where it counted.} For the round-3 changes we wrote the failing test first, watched it fail, implemented, and ran the whole suite. Some tests found real mistakes in our own designs before any experiment: the leakage test caught model names in packet source paths, and a unit test caught a risk-present rule that returned UNCERTAIN for a triggered integrity risk whenever an unrelated gate had also failed (a correction made before any packet was built, and recorded in the L2 pre-registration). The engineering issues in Table~\ref{tab:bugs} each have a regression test.

\paragraph{What the tests do not cover.} Most probe tests use injected fake backends, so they check logic and contracts, not behaviour on real checkpoints. The controlled experiments are the real-weight tests, and they found problems the unit tests could not (the label-mapping failure of Section~\ref{sec:pilot}). There is no test that the heuristic bands give sensible scores for an arbitrary degenerate model; the prediction-collapse gate covers the one case we met. The database-backed tests were not run in this session.

\section{Reproducibility}
\begin{itemize}
\item \textbf{Pinned inputs.} Civil Comments is pinned to a dataset revision; Hub models in the external audit are pinned to immutable revision SHAs that were re-resolved against the Hub before use; the selection file was written before any run; weight sha256 values are recorded.
\item \textbf{Seeds and determinism.} Probes take a seed; Fairness, Robustness and behavioural Safety evidence was identical across repeated runs in every experiment (zero standard deviation over five runs; L5 and L6 identical on repeats). A second data preparation with the same seed (L5) is byte-identical, including \code{ground\_truth.json}.
\item \textbf{Frozen, hashed results.} Result folders are never edited after the fact; the external audit folders carry a \code{SHA256SUMS} file; frozen v8 folders remain beside the v9 reruns; the L1 anchors are in a frozen file that the held-out run verifies against the code.
\item \textbf{Replayable analyses.} L2 stores every provider response keyed by packet hash and run index and can recompute all numbers with \code{--replay}; L3 can be replayed offline from stored inputs; the human analysis script re-reads the raw export (hash recorded) and regenerates the report. We re-ran the last of these while preparing this report and it reproduced the committed output exactly.
\item \textbf{Declared deviations.} Each plan document records post-run deviations with their timing (Table~\ref{tab:posthoc}).
\item \textbf{What is not reproducible.} LLM answers from live providers (quota, availability, and provider changes mean a re-run would differ); one run per Hub model; weights of the controlled models are not committed (they are reproducible from seeds, and their hashes are recorded).
\end{itemize}
Appendix~\ref{app:repro} lists the commands.

\section{Performance and resource use}
\begin{table}[H]
\centering\small
\caption{Measured cost on the reference laptop (RTX 4060, 8\,GB, CUDA, float32, batch size 8).}
\begin{tabularx}{\textwidth}{@{}p{5.6cm}L@{}}
\toprule
\textbf{Item} & \textbf{Measurement}\\
\midrule
Fine-tuning one BERT-base variant (8\,000 rows, 3 epochs) & 160.4--178.3\,s (recorded in the manifests)\\
One full evaluation of a real Hub model through the API (3\,000 rows, five probes, LLM call or fallback) & 23.2--110.5\,s over the seven v9 audit runs\\
One deterministic interpretation of a packet & about 2\,ms\\
One LLM call, successful & mean 19.8\,s, including in-run back-off before falling through to the next provider\\
Evaluation of a small model in the early benchmark (warm) & 5--15\,s wall-clock\\
Checkpoint size of BERT-base float32 & about 438\,MB\\
\bottomrule
\end{tabularx}
\end{table}
Everything ran sequentially, one model at a time, because the Docker stack shares about 16\,GB of RAM and 8\,GB of VRAM. TrustLens currently has no resource checker: it cannot tell before loading whether a model fits. We did not record peak memory per model, so we cannot say how close the largest models in our set (RoBERTa-large and XLM-R-large) came to the limit; a larger model would need a gate that does not exist yet.
